\section*{Chapter \ref{ch:g10:data-tables} --- Processing Data Tables}

\begin{solution}{exo:g10:data-tables:1}
More, or as many: every record that passed both tests still passes the weaker one, so the $4$ Hilltop days stay,
and Riverside's rainy days \omterm{def:g10:data-tables:join}{join} them. The result has $8$ records; it can never have fewer than $4$.
\end{solution}

\begin{solution}{exo:g10:data-tables:2}
\code{["fig", "pear", "kiwi", "apple"]}: the keys are the lengths $4, 3, 5, 4$. \code{"pear"} and \code{"kiwi"} have
the same key, $4$, and the sort is \omterm{def:g10:data-tables:sort-key}{stable}, so they keep their order.
\end{solution}

\begin{solution}{exo:g10:data-tables:3}
\code{return [{c: row[c] for c in cols} for row in table]} (\code{exercises/table_tools.py}): one new dictionary per
record, built in the order of \code{cols}.
\end{solution}

\begin{solution}{exo:g10:data-tables:4}
Key $1$: $2 \times 1 = 2$ pairs; key $2$: $1 \times 2 = 2$; key $3$ has no partner in $B$, key $4$ none in $A$. The
\omterm{def:g10:data-tables:join}{join} has $4$ records, and the keys $3$ and $4$ never appear in it.
\end{solution}

\begin{solution}{exo:g10:data-tables:5}
Without conversion the years are strings, and strings compare character by character: \code{"1"} comes before
\code{"9"}, so \code{"10"} comes before \code{"9"}. The key must be \code{int(r["year"])}
(\code{exercises/bugs.py}). For strings of the same length made of digits, the order of the strings is the order of
the numbers, which is why four-digit years hide the bug.
\end{solution}

\begin{solution}{exo:g10:data-tables:6}
Start from an empty dictionary and, for each record, do \code{counts[v] = counts.get(v, 0) + 1} with
\code{v = row[col]} (\code{exercises/table_tools.py}). On the books: science $13$, adventure $12$, history $11$,
comics $10$, mystery and sport $9$ each, poetry and nature $8$ each: $80$ in all.
\end{solution}

\begin{solution}{exo:g10:data-tables:7}
The two lines are different strings: one has a space that the other has not, so the set holds both. Compare records,
not lines: split each line at the commas, strip the spaces around each cell, and store the tuple of cells
(\code{exercises/bugs.py}).
\end{solution}

\begin{solution}{exo:g10:data-tables:8}
By class: (Ben, A), (Dee, A), (Ada, B), (Cy, B), (Eli, B); within a class the records keep their first order. By name
then by class: the names are already in order, so the first pass changes nothing and the result is the same, which
is also the table sorted by (class, name).
\end{solution}

\begin{solution}{exo:g10:data-tables:9}
Group the non-missing values in a dictionary of lists, then take the median of each list: sort it, and take the
middle value, or the mean of the two middle values when the length is even (\code{exercises/table_tools.py}). On the
raw Riverside record the median is $11.6$, against $11.65$ once cleaned; the mean is $11.43$, against $12.02$. Two
values of $-99$ move the median by one position among $363$, but they pull the mean by their whole size.
\end{solution}

\begin{solution}{exo:g10:data-tables:10}
$400 \times 80 = 32{,}000$ comparisons against $80 + 400 = 480$ steps. For $100{,}000$ loans and $20{,}000$ books:
$2 \times 10^9$ comparisons against $120{,}000$ steps.
\end{solution}

\begin{solution}{exo:g10:data-tables:11}
The two readings of $-99.0$ are numbers, not \code{None}: they are counted as temperatures and pull the mean down. The
corrected \code{mean(values, error_codes=(-99.0,))} keeps only the values that are neither \code{None} nor an error
code (\code{exercises/bugs.py}). With the duplicate removed as well, Riverside's mean is $12.02$.
\end{solution}

\begin{solution}{exo:g10:data-tables:12}
Build a dictionary from each key to the list of right records with that key. For each left record, add one merged
record per match, or, when there is none, the left record with \code{None} in every right column
(\code{exercises/table_tools.py}). On the town's tables nobody is left out: every member borrowed at least one book
and every book was lent at least once, so the left \omterm{def:g10:data-tables:join}{join} of the members with the loans has exactly the $400$ records of
the ordinary \omterm{def:g10:data-tables:join}{join}.
\end{solution}

\begin{solution}{exo:g10:data-tables:13}
Case 1: $k_1(r) < k_1(s)$. The second sort orders by $k_1$, so it puts $r$ before $s$, whatever came before. Case 2:
$k_1(r) = k_1(s)$ and $k_2(r) < k_2(s)$. After the first sort, $r$ is before $s$, since that sort orders by $k_2$. The
second sort sees equal keys $k_1$ for $r$ and $s$; being \omterm{def:g10:data-tables:sort-key}{stable}, it keeps them in that order. In both cases $r$ ends
before $s$, so the final table is sorted by $(k_1, k_2)$.
\end{solution}

\begin{solution}{exo:g10:data-tables:14}
For each loan, turn the dates into \code{date} objects; report the loan if its due date is not
\code{out + timedelta(days=21)}, or if a return date exists and is before \code{out} (\code{exercises/check_loans.py}). On the town's loans it finds no
problem; on a loan returned the day before it went out, with a due date $20$ days later, it reports both rules.
\end{solution}

\begin{solution}{exo:g10:data-tables:15}
\code{json.dumps} then \code{json.loads} gives back an equal list of dictionaries for the $60$ members and their
$400$ loans (\code{exercises/json_roundtrip.py}). Values that do not survive: a tuple comes back as a list, a
dictionary key that is a number comes back as a string, and sets, dates and bytes cannot be written at all without a
conversion chosen by the program.
\end{solution}

\begin{solution}{pb:g10:data-tables:1}
\textbf{1.} $2 \times 365 = 730$ records are expected; the table holds $731$.

\textbf{2.} Hilltop $365$ and Riverside $366$.

\textbf{3.} \code{date}: the days of 2025; \code{station}: Hilltop or Riverside; \code{temp_c}: temperatures, say
from $-30$ to $45$~°C; \code{rain_mm} and \code{wind_kmh}: numbers $0$ or more.

\textbf{4.} \code{None}.

\textbf{5.} Six missing temperatures, three at each station; two error codes, both at Riverside (29 January and
18 December); one duplicate, Riverside on 14 June.

\textbf{6.} \omterm{def:g10:data-tables:missing}{Missing values}: the test that a value is present. Error codes: the \omterm{def:g10:data-tables:header}{domain} of the temperature column.
Duplicate: the pair (date, station) must be unique.

\textbf{7.} Yes, every field is equal. It is most likely one reading saved twice, not two measurements.

\textbf{8.} Hilltop $362$ and Riverside $363$, the duplicate and the codes still included.

\textbf{9.} For the temperature mean, the record is left out; the rain and wind were measured and stay usable for
their own statistics. The program, which computes only temperatures, simply drops the record.

\textbf{10.} $0$~°C is a possible temperature: the mean would be wrong without any sign of it. The code means that
nothing was measured; it is a \omterm{def:g10:data-tables:missing}{missing value}.

\textbf{11.} The two are identical, so either; the program keeps the first one met.

\textbf{12.} It would assume that the two stations measure the same, or differ by a known amount; they differ by more
than a degree on average. Any filled value should be marked as estimated; here, leaving it missing is simpler and
honest.

\textbf{13.} $731 - 6 - 2 - 1 = 722$.

\textbf{14.} Hilltop: $10.628$~°C raw and corrected. Riverside: $11.430$ raw, $12.016$ corrected.

\textbf{15.} Hilltop's only defects are \omterm{def:g10:data-tables:missing}{missing values}, which the raw mean already skips; it has no code and no
duplicate.

\textbf{16.} $11.402$, which is $0.614$ below the corrected value: almost all the error comes from the codes; the
duplicate summer day raised the raw mean by about $0.03$.

\textbf{17.} The $360$ good values add up to about $360 \times 12.016 \approx 4325.8$. With the two codes:
$(4325.8 - 198) / 362 \approx 11.40$. Each code sits about $111$ degrees below the mean, and there are two among
$362$: $2 \times 111 / 362 \approx 0.61$.

\textbf{18.} July at both: $18.88$~°C at Hilltop and $21.25$~°C at Riverside, $2.37$ degrees warmer by the river.

\textbf{19.} The corrected 2025 mean temperature at Riverside is $12.02$~°C, over $360$ days; the two $-99.0$ error
codes would have lowered it to $11.40$~°C, an error of $0.61$ degree, and the raw file with its duplicate gives
$11.43$~°C.
\end{solution}
